\section{Метод решения}
Так как данная лабораторная работа по сути версией лабораторной работы 2, то был использован код из неё, дополненный функциями для file mapping'a (они реализованны только для систем POSIX). 
При запуске основной программы помимо родительского процесса запускаются два дочерних, все процессы связаны между собой через 3 mapped file'a, с системным вводом и выводом взаимодейтсвует только родительский процесс, ввод и вывод дочерних процессов перенаправляются в mapped file'ы

\section{Описание программы}
Общий интерфейс взаимодействия с операционной системой описан в заголовочном файле src/os\_interaction/basic/os.h.
При компиляции программы в ней выбирается та или иная реализация требуемых в работе функций, файлы с реализацией этих функций под разные системы находятся в папках src/os\_interaction/POSIX и src/os\_interaction/Windows
Основной код программы находится в файлах main.cpp, child1.cpp, child2.cpp

\textit{Осонвные структуры}
\begin{lstlisting}[language=C++]
struct FileHandleDummy { };
typedef FileHandleDummy* FileHandle;

struct ProcessHandleDummy { };
typedef ProcessHandleDummy* ProcessHandle;
\end{lstlisting}
Две структуры для описания HANDLE файлов и процессов для Windows.

\begin{lstlisting}[language=C++]
class Process {
public:
    int pid = -1;
    ProcessHandle handle = nullptr;
    Process(int pid_, ProcessHandle handle_): pid(pid_), handle(handle_) {};
    ~Process() {
        delete handle;
    }
};

class File {
public:
    int id = -1;
    FileHandle handle = nullptr;
    File(int id_, FileHandle handle_): id(id_), handle(handle_) {};
    File(std::string filename);
    ~File() {
        delete handle;
    }
};
\end{lstlisting}
Базовые классы для описания процессов в Windows и POSIX.

\begin{lstlisting}[language=C++]
struct ProcessStartInfo {

    std::string cmd;
    std::vector<std::string> args;

    File redirectStdin = {-1,nullptr};
    File redirectStdout = {-1,nullptr};

    ProcessStartInfo() = default;
};
\end{lstlisting}
Структура для описания параметров инициализации процесса (ключевые параметры - cmd, redirectStdin и redirectStdout).

\begin{lstlisting}[language=C++]
struct SharedMemory {
public:
    void* memptr = nullptr;
    File file;
    int memsize = 0;
    SharedMemory(void* ptr, File f, int memsize_): memptr(ptr), file(f), memsize(memsize_) {};
    ~SharedMemory();
};
\end{lstlisting}
Класс для описания file mapping'a, деструктор реализуется для каждой ОС отдельно.

\textit{Используемые системные вызовы}

Windows:
\begin{itemize}
  \item CreateProcess - создание процесса.
  \item TerminateProcess - принудительная остановка процесса.
  \item CreatePipe - создание pipe.
  \item GetStdHandle - получение стандартного потока ввода и вывода.
  \item Write - запись в поток.
  \item Read - чтение из потока.
\end{itemize}

POSIX:
\begin{itemize}
  \item posix\_spawn - создание процесса по сути является последовательным вызовом fork и execl (ещё можно между ними добавить dup2 или close), использован вместо fork и execl из-за похожести синтаксиса на CreateProcess, что повышает читаемость кода.
  \item kill - отправка процессу сигнала об окончании работы (SIGTERM или SIGKILL).
  \item pipe - создание pipe для POSIX.
  \item write - запись в поток.
  \item read - чтение из потока.
  \item flock - установка/удаление блокировки на файл для избежания гонок.
  \item mmap - создание shared memory и прикрепление её к file descriptor'у.
  \item munmap - освобождение shared memory.
  \item lseek - перемещение указателя для чтения из файла.
\end{itemize}
